% GENERATED FILE. Edit canonical lessons, then run npm run generate:books.
% canonical-source-hash: fnv1a64:c5a7aefc2a1aecff
% canonical-lessons: GU-C230-dadh, GU-C230-jadbu, GU-C230-asu, GU-C230-pidavu, GU-C230-bebhan-thavu, GU-R230-review-225-230

\chapter{Toothache and tears}
\label{ch:gu-toothache-and-tears}
\hlchaptermodality{230}

\begin{chapteropening}
\textbf{By the end of this chapter:} \emph{I can say that a tooth or my jaw hurts and ask whether someone fainted.}

Complete the last lesson of chapter 230: I can say that a tooth or my jaw hurts and ask whether someone fainted.
\end{chapteropening}

\section[\texorpdfstring{\gu{દાઢ} (dāḍh)}{દાઢ (dāḍh)}]{\gu{દાઢ} (dāḍh) — molar, back tooth}

\label{lesson:GU-C230-dadh}



\begin{quote}
Before the new one: say the Gujarati for heel, then the Gujarati for rib.
\end{quote}

\subsection*{You'll want to know: \gu{દાઢ}}
\textbf{\gu{દાઢ}} — \emph{dāḍh} — \textquotedblleft{}molar, back tooth\textquotedblright{}.

\begin{quote}
\textbf{\gu{મારી} \gu{દાઢ} \gu{હવે} \gu{નથી} \gu{દુખતી}.} — \emph{mārī dāḍh have nathī dukhtī} — \textquotedblleft{}My back tooth doesn't hurt any more.\textquotedblright{}
\end{quote}

Not \textbf{\gu{દાઢી}}, a beard.

\begin{cousinweb}
\textbf{\gu{દાઢ}} (\emph{dāḍh}) is a tadbhava. It comes from \emph{*dāṁṣṭra}, \textquotedblleft{}fang, tusk\textquotedblright{}, beside Sanskrit \emph{daṁṣṭra}, \textquotedblleft{}tusk, fang\textquotedblright{}; CDIAL glosses the Gujarati word \textquotedblleft{}molar\textquotedblright{} (CDIAL 6250).
\end{cousinweb}

\begin{grammarlens}[title={What you've built}]
One more word for everyday life at A2.
\end{grammarlens}

\subsection*{Your turn}
\begin{itemize}
\raggedright
  \item \emph{Say it:} \emph{dāḍh}
  \item \emph{Say it:} \emph{dāḍh}, once more
  \item \emph{Say it:} \emph{mārī dāḍh have nathī dukhtī} — My back tooth doesn't hurt any more.
  \item \emph{Recall:} say the Gujarati for heel, then the Gujarati for rib, then say \emph{dāḍh} again
\end{itemize}

\begin{tcolorbox}[breakable,colback=teal!4,colframe=teal!35!black,title={Before you move on}]
What does \gu{દાઢ} mean? (\textquotedblleft{}Molar, back tooth\textquotedblright{}.) Say it once more.
\end{tcolorbox}
\section[\texorpdfstring{\gu{જડબું} (jaḍbũ)}{જડબું (jaḍbũ)}]{\gu{જડબું} (jaḍbũ) — jaw}

\label{lesson:GU-C230-jadbu}



\begin{quote}
Before the new one: say the Gujarati for rib, then the Gujarati for molar.
\end{quote}

\subsection*{You'll want to know: \gu{જડબું}}
\textbf{\gu{જડબું}} — \emph{jaḍbũ} — \textquotedblleft{}jaw\textquotedblright{}.

\begin{quote}
\textbf{\gu{શું} \gu{તમારું} \gu{જડબું} \gu{દુખે} \gu{છે}?} — \emph{śũ tamārũ jaḍbũ dukhe che?} — \textquotedblleft{}Does your jaw hurt?\textquotedblright{}
\end{quote}

\begin{cousinweb}
The origin of \textbf{\gu{જડબું}} (\emph{jaḍbũ}) is not settled: CDIAL lists no Gujarati form for it, and none of its entries for the jaw holds the word.
\end{cousinweb}

\begin{grammarlens}[title={What you've built}]
One more word for everyday life at A2.
\end{grammarlens}

\subsection*{Your turn}
\begin{itemize}
\raggedright
  \item \emph{Say it:} \emph{jaḍbũ}
  \item \emph{Say it:} \emph{jaḍbũ}, once more
  \item \emph{Say it:} \emph{śũ tamārũ jaḍbũ dukhe che?} — Does your jaw hurt?
  \item \emph{Recall:} say the Gujarati for rib, then the Gujarati for molar, then say \emph{jaḍbũ} again
\end{itemize}

\begin{tcolorbox}[breakable,colback=teal!4,colframe=teal!35!black,title={Before you move on}]
What does \gu{જડબું} mean? (\textquotedblleft{}Jaw\textquotedblright{}.) Say it once more.
\end{tcolorbox}
\section[\texorpdfstring{\gu{આંસુ} (ā̃su)}{આંસુ (ā̃su)}]{\gu{આંસુ} (ā̃su) — tear, teardrop}

\label{lesson:GU-C230-asu}



\begin{quote}
Before the new one: say the Gujarati for molar, then the Gujarati for jaw.
\end{quote}

\subsection*{You'll want to know: \gu{આંસુ}}
\textbf{\gu{આંસુ}} — \emph{ā̃su} — \textquotedblleft{}tear, teardrop\textquotedblright{}.

\begin{quote}
\textbf{\gu{તેની} \gu{આંખમાં} \gu{આંસુ} \gu{નથી}.} — \emph{tenī ā̃khmā̃ ā̃su nathī} — \textquotedblleft{}There are no tears in her eyes.\textquotedblright{}
\end{quote}

A neuter noun.

\begin{cousinweb}
\textbf{\gu{આંસુ}} (\emph{ā̃su}) is a tadbhava, a word inherited from Sanskrit and reshaped by sound change. It comes from \emph{aśru}, \textquotedblleft{}tear\textquotedblright{}, which picked up its nasal early, as the word for eye did (CDIAL 919).
\end{cousinweb}

\begin{grammarlens}[title={What you've built}]
One more word for everyday life at A2.
\end{grammarlens}

\subsection*{Your turn}
\begin{itemize}
\raggedright
  \item \emph{Say it:} \emph{ā̃su}
  \item \emph{Say it:} \emph{ā̃su}, once more
  \item \emph{Say it:} \emph{tenī ā̃khmā̃ ā̃su nathī} — There are no tears in her eyes.
  \item \emph{Recall:} say the Gujarati for molar, then the Gujarati for jaw, then say \emph{ā̃su} again
\end{itemize}

\begin{tcolorbox}[breakable,colback=teal!4,colframe=teal!35!black,title={Before you move on}]
What does \gu{આંસુ} mean? (\textquotedblleft{}Tear, teardrop\textquotedblright{}.) Say it once more.
\end{tcolorbox}
\section[\texorpdfstring{\gu{પીડાવું} (pīḍāvũ)}{પીડાવું (pīḍāvũ)}]{\gu{પીડાવું} (pīḍāvũ) — to suffer, to be in pain}

\label{lesson:GU-C230-pidavu}



\begin{quote}
Before the new one: say the Gujarati for jaw, then the Gujarati for tear.
\end{quote}

\subsection*{You'll want to know: \gu{પીડાવું}}
\textbf{\gu{પીડાવું}} — \emph{pīḍāvũ} — \textquotedblleft{}to suffer, to be in pain\textquotedblright{}.

\begin{quote}
\textbf{\gu{શું} \gu{દર્દી} \gu{પીડાય} \gu{છે}?} — \emph{śũ dardī pīḍāy che?} — \textquotedblleft{}Is the patient in pain?\textquotedblright{}
\end{quote}

\begin{cousinweb}
\textbf{\gu{પીડાવું}} (\emph{pīḍāvũ}) is built on \emph{pīḍā}, \textquotedblleft{}pain, damage\textquotedblright{}, a Sanskrit word taken over whole; CDIAL lists the inherited Gujarati \emph{pīḍ}, \textquotedblleft{}pain\textquotedblright{}, under it (CDIAL 8227).
\end{cousinweb}

\begin{grammarlens}[title={What you've built}]
One more word for everyday life at A2.
\end{grammarlens}

\subsection*{Your turn}
\begin{itemize}
\raggedright
  \item \emph{Say it:} \emph{pīḍāvũ}
  \item \emph{Say it:} \emph{pīḍāvũ}, once more
  \item \emph{Say it:} \emph{śũ dardī pīḍāy che?} — Is the patient in pain?
  \item \emph{Recall:} say the Gujarati for jaw, then the Gujarati for tear, then say \emph{pīḍāvũ} again
\end{itemize}

\begin{tcolorbox}[breakable,colback=teal!4,colframe=teal!35!black,title={Before you move on}]
What does \gu{પીડાવું} mean? (\textquotedblleft{}To suffer, to be in pain\textquotedblright{}.) Say it once more.
\end{tcolorbox}
\section[\texorpdfstring{\gu{બેભાન} \gu{થવું} (bebhān thavũ)}{બેભાન થવું (bebhān thavũ)}]{\gu{બેભાન} \gu{થવું} (bebhān thavũ) — to faint, to pass out}

\label{lesson:GU-C230-bebhan-thavu}



\begin{quote}
Before the new one: say the Gujarati for tear, then the Gujarati for to suffer.
\end{quote}

\subsection*{You'll want to know: \gu{બેભાન} \gu{થવું}}
\textbf{\gu{બેભાન} \gu{થવું}} — \emph{bebhān thavũ} — \textquotedblleft{}to faint, to pass out\textquotedblright{}.

\begin{quote}
\textbf{\gu{શું} \gu{તે} \gu{બેભાન} \gu{થયો}?} — \emph{śũ te bebhān thayo?} — \textquotedblleft{}Did he faint?\textquotedblright{}
\end{quote}

Word for word, to become senseless: \textbf{\gu{બેભાન}}, \emph{bebhān}, and \textbf{\gu{થવું}}, \emph{thavũ}.

\begin{cousinweb}
\textbf{\gu{બેભાન} \gu{થવું}} (\emph{bebhān thavũ}) is built from \emph{bebhān}, \textquotedblleft{}senseless\textquotedblright{}, and \emph{thavũ}, \textquotedblleft{}to become\textquotedblright{}. \emph{Bebhān} is usually explained as the Persian prefix \emph{be-}, \textquotedblleft{}without\textquotedblright{}, joined to \emph{bhān}, \textquotedblleft{}awareness\textquotedblright{}, usually traced to Sanskrit \emph{bhāna}, \textquotedblleft{}appearance\textquotedblright{} (CDIAL 9451); CDIAL lists no Gujarati form of the whole.
\end{cousinweb}

\begin{grammarlens}[title={What you've built}]
One more word for everyday life at A2.
\end{grammarlens}

\subsection*{Your turn}
\begin{itemize}
\raggedright
  \item \emph{Say it:} \emph{bebhān thavũ}
  \item \emph{Say it:} \emph{bebhān thavũ}, once more
  \item \emph{Say it:} \emph{śũ te bebhān thayo?} — Did he faint?
  \item \emph{Recall:} say the Gujarati for tear, then the Gujarati for to suffer, then say \emph{bebhān thavũ} again
\end{itemize}

\begin{tcolorbox}[breakable,colback=teal!4,colframe=teal!35!black,title={Before you move on}]
What does \gu{બેભાન} \gu{થવું} mean? (\textquotedblleft{}To faint, to pass out\textquotedblright{}.) Say it once more.
\end{tcolorbox}
\section[(dialogue)]{Words from chapters 225–230 — a review}

\label{lesson:GU-R230-review-225-230}



\begin{quote}
Say the Gujarati for dizziness and to vomit.
\end{quote}

\subsection*{Your turn}
Say these aloud:

\begin{itemize}
\raggedright
  \item dizziness, to vomit, to digest, to yawn, sweat
  \item to shiver, to pant, to sniff, to swallow, lung
  \item injury, wound, bone, to sprain, to limp
  \item swelling, to swell, to scratch, skin, blister
  \item wrist, thumb, ankle, heel, rib
  \item molar, jaw, tear, to suffer, to faint
\end{itemize}

\begin{tcolorbox}[breakable,colback=teal!4,colframe=teal!35!black,title={Before you move on}]
Dizziness? (\textbf{\gu{ચક્કર}}.) To vomit? (\textbf{\gu{ઊલટી} \gu{થવી}}.) To digest? (\textbf{\gu{પચવું}}.) To yawn? (\textbf{\gu{બગાસું} \gu{ખાવું}}.) Sweat? (\textbf{\gu{પરસેવો}}.) To shiver? (\textbf{\gu{ધ્રૂજવું}}.) To pant? (\textbf{\gu{હાંફવું}}.) To sniff? (\textbf{\gu{સૂંઘવું}}.) To swallow? (\textbf{\gu{ગળવું}}.) Lung? (\textbf{\gu{ફેફસું}}.) Injury? (\textbf{\gu{ઈજા}}.) Wound? (\textbf{\gu{ઘા}}.) Bone? (\textbf{\gu{હાડકું}}.) To sprain? (\textbf{\gu{મચકોડાવું}}.) To limp? (\textbf{\gu{લંગડાવું}}.) Swelling? (\textbf{\gu{સોજો}}.) To swell? (\textbf{\gu{સૂજવું}}.) To scratch? (\textbf{\gu{ખંજવાળવું}}.) Skin? (\textbf{\gu{ચામડી}}.) Blister? (\textbf{\gu{ફોલ્લો}}.) Wrist? (\textbf{\gu{કાંડું}}.) Thumb? (\textbf{\gu{અંગૂઠો}}.) Ankle? (\textbf{\gu{ઘૂંટી}}.) Heel? (\textbf{\gu{એડી}}.) Rib? (\textbf{\gu{પાંસળી}}.) Molar? (\textbf{\gu{દાઢ}}.) Jaw? (\textbf{\gu{જડબું}}.) Tear? (\textbf{\gu{આંસુ}}.) To suffer? (\textbf{\gu{પીડાવું}}.) To faint? (\textbf{\gu{બેભાન} \gu{થવું}}.)
\end{tcolorbox}
